% !TeX program = xelatex
\documentclass[11pt,a4paper]{ctexart}

\usepackage[left=2.4cm,right=2.4cm,top=2.1cm,bottom=2.1cm]{geometry}
\usepackage{amsmath,amssymb,mathtools,braket}
\usepackage{xcolor}
\usepackage[inline]{enumitem}
\usepackage{fancyhdr}
\usepackage[hidelinks]{hyperref}

% ============================================================
% 基本信息
% ============================================================
\newcommand{\coursename}{高等量子力学}
\newcommand{\assignmentnumber}{第 2 次作业}
\newcommand{\teachername}{\underline{\makebox[2.8cm][c]{金彪}}}
\newcommand{\studentname}{\underline{\makebox[2.8cm][c]{林照翔}}}
\newcommand{\studentid}{\underline{\makebox[3.0cm][c]{202628001207046}}}

% ============================================================
% 页面样式
% ============================================================
\setlength{\parindent}{2em}
\setlength{\parskip}{0.25em}
\setlength{\headheight}{14pt}
\pagestyle{fancy}
\fancyhf{}
\lhead{\coursename}
\rhead{\assignmentnumber}
\cfoot{\thepage}
\renewcommand{\headrulewidth}{0.4pt}

% ============================================================
% 列表参数
% ============================================================
\setlist[enumerate,1]{
  label=(\alph*),
  leftmargin=3.2em,
  itemsep=0.55em,
  topsep=0.45em
}

\newlist{inlineenum}{enumerate*}{1}
\setlist[inlineenum,1]{
  label=(\alph*),
  labelsep=0.5em,
  itemjoin={；\quad},
  itemjoin*={；\quad}
}

% ============================================================
% 题目与回答空间
% ============================================================
\newcounter{problem}
\newenvironment{problem}
  {\refstepcounter{problem}%
   \par\medskip\noindent
   \begingroup
   \setlength{\fboxsep}{3pt}%
   \setlength{\fboxrule}{0.8pt}%
   \fbox{\bfseries 题目\chinese{problem}}%
   \endgroup
   \quad\ignorespaces}
  {\par\medskip}

\newenvironment{answer}
  {\par\smallskip\noindent
   \begingroup
   \setlength{\fboxsep}{4pt}%
   \colorbox{black}{\textcolor{white}{\bfseries 回答}}%
   \endgroup
   \quad\ignorespaces}
  {\par\smallskip\noindent\rule{\linewidth}{0.4pt}\medskip}

\newcommand{\answerspace}[1]{\par\vspace*{#1}}

\begin{document}

\begin{center}
  {\LARGE\bfseries \coursename\quad \assignmentnumber}\par
  \vspace{0.7em}
  \small
  教师：\teachername\qquad
  姓名：\studentname\qquad
  学号：\studentid
\end{center}
\vspace{0.4em}

% 本份作业承接原题前五题，题号从六开始。
\setcounter{problem}{5}

\begin{problem}
自旋为 $1/2$ 的粒子，其三个自旋分量算符在 $s_z$ 表象下的矩阵表示分别为
\begin{equation}
  \hat s_x=\frac{\hbar}{2}
  \left(\begin{matrix}0&1\\1&0\end{matrix}\right),\qquad
  \hat s_y=\frac{\hbar}{2}
  \left(\begin{matrix}0&-\mathrm{i}\\\mathrm{i}&0\end{matrix}\right),\qquad
  \hat s_z=\frac{\hbar}{2}
  \left(\begin{matrix}1&0\\0&-1\end{matrix}\right).
\end{equation}
根据幺正变换矩阵的定义式：
\begin{enumerate}
  \item 求由 $s_z$ 表象到 $s_y$ 表象的幺正变换矩阵 $S_{zy}$；
  \item 求由 $s_z$ 表象到 $s_x$ 表象的幺正变换矩阵 $S_{zx}$。
\end{enumerate}
\end{problem}

\begin{answer}
  \begin{enumerate}
    \item 对于两个力学量算符 $\hat{A},\hat{B} $，设它们的本征方程分别为 
      \begin{equation}
        \hat{A} \Ket{a_i } = a_i \Ket{a_i },\qquad
        \hat{B} \Ket{b_i } = b_i \Ket{b_i },
      \end{equation}
    则从 $A $ 表象到 $B $ 表象的幺正变换算符 $\hat{S}_{A\to B} $ 满足
      \begin{equation}
        \hat{S}_{A\to B} = \sum_{i} \Ket{b_i } \Bra{a_i },\qquad
        \hat{S}_{A\to B} \Ket{a_i } = \Ket{b_i },\qquad
        \left(S_{A\to B} \right)_{ij} = \Braket{a_i | b_j }.
      \end{equation}
      因此，$s_z$ 表象到 $s_y$ 表象的幺正变换矩阵 $S_{z\to y}$ 形式上可写为
      \begin{equation}
        S_{z\to y}
        =\begin{pmatrix}
        \Braket{+_z|+_y } &\Braket{+_z|-_y } \\
        \Braket{-_z|+_y } &\Braket{-_z|-_y }
        \end{pmatrix}.
      \end{equation}
      $s_z $ 表象下，$\hat{s}_z,\hat{s}_y $ 本征态的矩阵表示为
      \begin{equation}
        \Ket{+_z }
        =\begin{pmatrix}
        1 \\
        0
        \end{pmatrix},\qquad
        \Ket{-_z }
        =\begin{pmatrix}
        0 \\
        1
        \end{pmatrix},\qquad
        \Ket{+_y }
        =\frac{1 }{\sqrt{2} } 
        \begin{pmatrix}
          1 \\
          \mathrm{i}
        \end{pmatrix},\qquad
        \Ket{-_y }
        =\frac{1 }{\sqrt{2} }
        \begin{pmatrix}
          -1 \\
          \mathrm{i}
        \end{pmatrix}.
      \end{equation}
      内积为
      \begin{equation}
        \Braket{+_z | +_y } = \frac{1 }{\sqrt{2} } ,\qquad
        \Braket{+_z | -_y } = -\frac{1 }{\sqrt{2} } ,\qquad
        \Braket{-_z | +_y } = \frac{\mathrm{i} }{\sqrt{2} } ,\qquad
        \Braket{-_z | -_y } = \frac{\mathrm{i} }{\sqrt{2} }, 
      \end{equation}
      于是 $S_{z\to y} $ 为
      \begin{equation}
        S_{z\to y}
        =\begin{pmatrix}
          \frac{1 }{\sqrt{2} } &-\frac{1 }{\sqrt{2} } \\
          \frac{\mathrm{i} }{\sqrt{2} } &\frac{\mathrm{i} }{\sqrt{2} }
        \end{pmatrix}.
      \end{equation}
    \item $s_z $ 表象下，$\hat{s}_x $ 本征态的矩阵表示为
    \begin{equation}
      \Ket{+_x }
      =\frac{1 }{\sqrt{2} }
      \begin{pmatrix}
        1 \\
        1
      \end{pmatrix},\qquad
      \Ket{-_x }
      =\frac{1 }{\sqrt{2} } 
      \begin{pmatrix}
        1 \\
        -1
      \end{pmatrix}.
    \end{equation}
    于是由 $s_z$ 表象到 $s_x$ 表象的幺正变换矩阵 $S_{z\to x}$ 为
    \begin{equation}
      S_{z\to x}
      =\begin{pmatrix}
      \Braket{+_z|+_x } &\Braket{+_z|-_x } \\
      \Braket{-_z|+_x } &\Braket{-_z|-_x }
      \end{pmatrix}
      =\begin{pmatrix}
        \frac{1 }{\sqrt{2} } &\frac{1 }{\sqrt{2} } \\
        \frac{1 }{\sqrt{2} } &\frac{-1 }{\sqrt{2} } 
      \end{pmatrix}.
    \end{equation}
  \end{enumerate}
\end{answer}

\begin{problem}
从 $\hat a\Ket{0}=0$ 出发，求谐振子动量波函数中基态 $\Psi_{p0}$ 及第一激发态
$\Psi_{p1}$ 的具体表达式，其中 $\Psi_{pn}=\Braket{p|n}$。
\end{problem}

\begin{answer}
  湮灭算符可用坐标、动量算符表达：
  \begin{equation}
    \hat{a} = \sqrt{\frac{m\omega }{2\hbar } } \left(\hat{x} + \frac{\mathrm{i} }{m\omega } \hat{p} \right).
  \end{equation}
  $\hat{a} \Ket{0 } = 0 $ 左乘 $\Bra{p } $ 有：
  \begin{equation}
    \Braket{p | \hat{x} | 0 } + \frac{\mathrm{i} }{m\omega } \Braket{p | \hat{p} | 0 }
    =0.
  \end{equation}
  注意到：
  \begin{equation}
    \Braket{p | \hat{p} | 0 } = p \Braket{p | 0 } = p \Psi_{p0},
  \end{equation}
  \begin{equation}
    \begin{aligned}
      \Braket{p | \hat{x} | 0 }
      &=\int \mathrm{d}p' \Braket{p|\hat{x}|p' } \Braket{p'|0 }
      =\int \mathrm{d}p' \Braket{p|\hat{x}|p' } \Psi_{p'0},
    \end{aligned}
  \end{equation}
  其中
  \begin{equation}
    \begin{aligned}
      \Braket{p|\hat{x}|p' }
      &=\int \mathrm{d}x \mathrm{d}x' \Braket{p|x }\Braket{x|\hat{x}|x' } \Braket{x'|p' }
      =\frac{1 }{2\pi \hbar } \int \mathrm{d}x\mathrm{d}x' x \mathrm{e}^{\mathrm{i}\left(p'x'-px \right)/\hbar} \delta(x-x') \\
      &=\frac{1 }{2\pi\hbar } \int \mathrm{d}x x\mathrm{e}^{\mathrm{i}\left(p'-p \right)x/\hbar}
      =\frac{1 }{2\pi\hbar } \frac{\hbar }{\mathrm{i} }  \int \mathrm{d}x \frac{\partial }{\partial \left(p'-p \right) } \mathrm{e}^{\mathrm{i}\left(p'-p \right)x/\hbar} \\
      &=\frac{1 }{2\pi \mathrm{i} } \frac{\mathrm{d} }{\mathrm{d}\left(p'-p \right) } \int \mathrm{d}x \mathrm{e}^{\mathrm{i}\left(p'-p \right)x/\hbar}
      =\frac{1 }{\mathrm{i} } \frac{\mathrm{d} }{\mathrm{d}\left(p'-p \right) } \delta\left(\frac{p'-p }{\hbar }  \right) \\
      &=-\mathrm{i}\hbar \frac{\mathrm{d}\delta(p'-p) }{\mathrm{d}\left(p'-p \right) }
      =-\mathrm{i}\hbar \delta'(p'-p).
    \end{aligned}
  \end{equation}
  于是
  \begin{equation}
    \begin{aligned}
      \Braket{p | \hat{x} | 0 }
      &=\int \mathrm{d}p' \Braket{p|\hat{x}|p' } \Psi_{p'0}
      =-\mathrm{i} \hbar \int \mathrm{d}p' \delta'(p'-p) \Psi_{p'0} \\
      &=\mathrm{i}\hbar \frac{\mathrm{d} }{\mathrm{d}p } \Psi_{p0}.
    \end{aligned}
  \end{equation}
  因此，$\Psi_{p0} $ 满足微分方程
  \begin{equation}
    \mathrm{i}\hbar \frac{\mathrm{d} }{\mathrm{d}p } \Psi_{p0} + \frac{\mathrm{i} }{m\omega } p \Psi_{p0} = 0, 
  \end{equation}
  分离变量得：
  \begin{equation}
    \frac{\mathrm{d}\Psi_{p0} }{\Psi_{p0} } = -\frac{1 }{m\omega \hbar } p \mathrm{d}p.
  \end{equation}
  两边积分得到通解
  \begin{equation}
    \Psi_{p0} = C_0 \exp\left(-\frac{p^2 }{2m\omega\hbar }  \right).
  \end{equation}
  利用高斯积分公式
  \begin{equation}
    \int_{-\infty}^{+\infty} \mathrm{e}^{-ax^2}\mathrm{d}x
    =\sqrt{\frac{\pi }{a } },
  \end{equation}
  利用归一化条件
  \begin{equation}
    \int_{-\infty}^{+\infty} \left|\Psi_{p0} \right|^2\mathrm{d}p=1,
  \end{equation}
  归一化系数可取为
  \begin{equation}
    C_0 = \left(\pi m\omega\hbar \right)^{-1/4},
  \end{equation}
  于是
  \begin{equation}
    \Psi_{p0} = \left(\pi m\omega\hbar \right)^{-1/4} \exp\left(-\frac{p^2 }{2m\omega\hbar } \right).
  \end{equation}
  产生算符同样可有坐标、动量算符表示：
  \begin{equation}
    \hat{a}^\dag
    =\sqrt{\frac{m\omega }{2\hbar } }\left(\hat{x}-\frac{\mathrm{i} }{m\omega } \hat{p} \right).
  \end{equation}
  于是
  \begin{equation}
    \begin{aligned}
      \Psi_{p1}
      &=\Braket{p|1 }
      =\Braket{p|\hat{a}^\dag|0 }
      =\sqrt{\frac{m\omega }{2\hbar } }\left(\Braket{p|\hat{x}|0 } - \frac{\mathrm{i} }{m\omega }\Braket{p|\hat{p}|0 }  \right) \\
      &=\sqrt{\frac{m\omega }{2\hbar } } \left(\mathrm{i}\hbar\frac{\mathrm{d} }{\mathrm{d}p } \Psi_{p0} - \frac{\mathrm{i} }{m\omega } p\Psi_{p0}  \right)
      =\sqrt{\frac{m\omega }{2\hbar } } \frac{\left(-2\mathrm{i} \right) }{m\omega } p \Psi_{p0}
      =-\mathrm{i}\sqrt{\frac{2 }{m\omega\hbar } } p \Psi_{p0} \\
      &=-\mathrm{i}\sqrt{\frac{2 }{m\omega\hbar } } \left(\pi m\omega\hbar \right)^{-1/4} p \exp\left(-\frac{p^2 }{2m\omega\hbar } \right).
    \end{aligned}
  \end{equation}
\end{answer}

\begin{problem}
\begin{enumerate}
  \item 求证谐振子系统坐标和动量算符本征态的下列展开式：
  \begin{align}
    \Ket{x}
    &=\sum_{n=0}^{\infty}\Ket{n}\Braket{n|x}\notag\\
    &=\left(\frac{m\omega}{\pi\hbar}\right)^{1/4}\exp\left[-\frac{m\omega}{2\hbar}x^2+\sqrt{\frac{2m\omega}{\hbar}}x\hat a^\dagger-\frac{\left(\hat a^\dagger\right)^2}{2}\right]\Ket{0},\\
    \Ket{p}
    &=\sum_{n=0}^{\infty}\Ket{n}\Braket{n|p}\notag\\
    &=\left(\frac{1}{m\omega\pi\hbar}\right)^{1/4}\exp\left[-\frac{p^2}{2m\omega\hbar}+\sqrt{\frac{2}{m\omega\hbar}}\,\mathrm{i}p\hat a^\dagger+\frac{\left(\hat a^\dagger\right)^2}{2}\right]\Ket{0}.
  \end{align}
  \item 验证上述表达式满足
  \begin{equation}
    \hat x\Ket{x}=x\Ket{x},\qquad
    \hat p\Ket{p}=p\Ket{p}.
  \end{equation}
\end{enumerate}
\end{problem}

\begin{answer}
  \begin{enumerate}
    \item 谐振子本征态坐标表象波函数为
    \begin{equation}
      \psi_n(x)
      =\Braket{x|n }
      =\frac{1 }{\pi^{1/4} \sqrt{x_0} } \frac{H_n(\xi) }{\sqrt{2^n n!} } \mathrm{e}^{-\xi^2/2},\quad 
      x_0 = \sqrt{\frac{\hbar }{m\omega } },\quad
      \xi = \frac{x }{x_0 } = \sqrt{\frac{m\omega }{\hbar } } x.
    \end{equation}
    利用厄米多项式的生成函数公式
    \begin{equation}
      \mathrm{e}^{-t^2+2xt} = \sum_{n=0}^{\infty} \frac{H_n(x) }{n! } t^n,
    \end{equation}
    可得
    \begin{equation}
      \begin{aligned}
        \Ket{x }
        &=\sum_{n=0}^{\infty} \Braket{n|x } \Ket{n }
        =\sum_{n=0}^{\infty} \psi_n^*(x) \frac{\left(\hat{a}^\dag \right)^n }{\sqrt{n!} } \Ket{0 }
        =\frac{1 }{\pi^{1/4} \sqrt{x_0} } \mathrm{e}^{-\xi^2/2} \left[\sum_{n=0}^{\infty} \frac{H_n(\xi) }{n! } \left(\frac{\hat{a}^\dag }{\sqrt{2} }  \right)^n  \right] \Ket{0 } \\
        &=\frac{1 }{\pi^{1/4} \sqrt{x_0} } \mathrm{e}^{-\xi^2/2} \exp\left[-\frac{\left(\hat{a}^\dag \right)^2  }{2 } + \sqrt{2} \xi \hat{a}^\dag \right] \Ket{0 } \\
        &=\frac{1 }{\pi^{1/4} \sqrt{x_0} } \exp\left[-\frac{\left(\hat{a}^\dag \right)^2  }{2 } + \sqrt{2} \xi \hat{a}^\dag - \frac{\xi^2 }{2 }  \right] \Ket{0 } \\
        &=\left(\frac{m\omega}{\pi\hbar}\right)^{1/4}\exp\left[-\frac{m\omega}{2\hbar}x^2+\sqrt{\frac{2m\omega}{\hbar}}x\hat a^\dagger-\frac{\left(\hat a^\dagger\right)^2}{2}\right]\Ket{0}.
      \end{aligned}
    \end{equation}
    利用高斯积分公式
    \begin{equation}
      \int_{-\infty}^{+\infty} \mathrm{e}^{-Ax^2+Bx} \mathrm{d}x
      =\sqrt{\frac{\pi }{A } } \exp\left(\frac{B^2 }{4A }  \right),
    \end{equation}
    可得
    \begin{equation}
      \begin{aligned}
        \Ket{p }
        &=\int \mathrm{d}x \Braket{x|p } \Ket{x } \\
        &=\left(\frac{m\omega}{\pi\hbar}\right)^{1/4} \frac{1 }{\sqrt{2\pi\hbar} } \int\mathrm{d}x \mathrm{e}^{\mathrm{i}px/\hbar} \exp\left[-\frac{m\omega}{2\hbar}x^2+\sqrt{\frac{2m\omega}{\hbar}}x\hat a^\dagger-\frac{\left(\hat a^\dagger\right)^2}{2}\right] \Ket{0 } \\
        &=\left(\frac{m\omega}{\pi\hbar}\right)^{1/4} \frac{1 }{\sqrt{2\pi\hbar} } \int\mathrm{d}x \exp\left[-\frac{m\omega}{2\hbar}x^2+\left(\sqrt{\frac{2m\omega}{\hbar}}\hat a^\dagger + \frac{\mathrm{i}p }{\hbar } \right)x  -\frac{\left(\hat a^\dagger\right)^2}{2}\right] \Ket{0 } \\
        &=\left(\frac{m\omega}{\pi\hbar}\right)^{1/4} \frac{1 }{\sqrt{2\pi\hbar} } \sqrt{\frac{2\pi\hbar }{m\omega } } \exp\left\{ \frac{\hbar }{2m\omega } \left[\frac{2m\omega }{\hbar } \left(\hat{a}^\dag \right)^2 + \frac{2\mathrm{i} }{\hbar } \sqrt{\frac{2m\omega }{\hbar } } p \hat{a}^\dag - \frac{p^2 }{\hbar^2 }   \right] \right\} \exp\left[-\frac{\left(\hat{a}^\dag \right)^2 }{2 } \right] \Ket{0 } \\
        &=\left(\frac{1}{m\omega\pi\hbar}\right)^{1/4}\exp\left[-\frac{p^2}{2m\omega\hbar}+\sqrt{\frac{2}{m\omega\hbar}}\,\mathrm{i}p\hat a^\dagger+\frac{\left(\hat a^\dagger\right)^2}{2}\right]\Ket{0}.
      \end{aligned}
    \end{equation}
    \item 由
    \begin{equation}
      \hat{a}
      =\sqrt{\frac{m\omega }{2\hbar } } \left(\hat{x} + \frac{\mathrm{i} }{m\omega } \hat{p} \right),\qquad
      \hat{a}^\dag
      =\sqrt{\frac{m\omega }{2\hbar } } \left(\hat{x} - \frac{\mathrm{i} }{m\omega } \hat{p} \right),
    \end{equation}
    可以反解出
    \begin{equation}
      \hat{x}
      =\sqrt{\frac{\hbar }{2m\omega } } \left(\hat{a}+\hat{a}^\dag \right),\qquad
      \hat{p}
      =-\mathrm{i} \sqrt{\frac{\hbar m \omega }{2 } } \left(\hat{a}-\hat{a}^\dag \right).
    \end{equation}
    利用
    \begin{equation}
      \hat{a} f\left(\hat{a}^\dag \right) \Ket{0 }
      =\left[\hat{a},f\left(\hat{a}^\dag \right) \right] \Ket{0 },\qquad
      \left[\hat{a},f\left(\hat{a}^\dag \right) \right]
      =f'\left(\hat{a}^\dag \right),
    \end{equation}
    令
    \begin{equation}
      F\left(\hat{a}^\dag \right)
      =\exp\left[-\frac{m\omega}{2\hbar}x^2+\sqrt{\frac{2m\omega}{\hbar}}x\hat a^\dagger-\frac{\left(\hat a^\dagger\right)^2}{2}\right],
    \end{equation}
    \begin{equation}
      \begin{aligned}
        F'\left(\hat{a}^\dag \right)
      &=\left(\sqrt{\frac{2m\omega }{\hbar } }x - \hat{a}^\dag \right)\exp\left[-\frac{m\omega}{2\hbar}x^2+\sqrt{\frac{2m\omega}{\hbar}}x\hat a^\dagger-\frac{\left(\hat a^\dagger\right)^2}{2}\right] \\
      &=\left(\sqrt{\frac{2m\omega }{\hbar } }x - \hat{a}^\dag \right) F\left(\hat{a}^\dag \right)
      \end{aligned}
    \end{equation}
    可得
    \begin{equation}
      \begin{aligned}
        \hat{x} \Ket{x }
        &=\sqrt{\frac{\hbar }{2m\omega } } \left(\hat{a}+\hat{a}^\dag \right) \left\{\left(\frac{m\omega}{\pi\hbar}\right)^{1/4}\exp\left[-\frac{m\omega}{2\hbar}x^2+\sqrt{\frac{2m\omega}{\hbar}}x\hat a^\dagger-\frac{\left(\hat a^\dagger\right)^2}{2}\right]\Ket{0} \right\} \\
        &=\sqrt{\frac{\hbar }{2m\omega } }\left(\frac{m\omega}{\pi\hbar}\right)^{1/4} \left(\hat{a}+\hat{a}^\dag \right) F\left(\hat{a}^\dag \right) \Ket{0 } \\
        &=\sqrt{\frac{\hbar }{2m\omega } }\left(\frac{m\omega}{\pi\hbar}\right)^{1/4} \left\{\left[\hat{a},F\left(\hat{a}^\dag \right) \right]+\hat{a}^\dag F\left(\hat{a}^\dag \right) \right\} \Ket{0 } \\
        &=\sqrt{\frac{\hbar }{2m\omega } }\left(\frac{m\omega}{\pi\hbar}\right)^{1/4} \left\{F'\left(\hat{a}^\dag \right)+\hat{a}^\dag F\left(\hat{a}^\dag \right) \right\} \Ket{0 } \\
        &=\sqrt{\frac{\hbar }{2m\omega } }\left(\frac{m\omega}{\pi\hbar}\right)^{1/4} \left\{\left(\sqrt{\frac{2m\omega }{\hbar } }x - \hat{a}^\dag \right) F\left(\hat{a}^\dag \right)+\hat{a}^\dag F\left(\hat{a}^\dag \right) \right\} \Ket{0 } \\
        &=x \left(\frac{m\omega}{\pi\hbar}\right)^{1/4} F\left(\hat{a}^\dag \right) \Ket{0 } \\
        &=x\Ket{x }.
      \end{aligned}
    \end{equation}
    令
    \begin{equation}
      G\left(\hat{a}^\dag \right)
      =\exp\left[-\frac{p^2}{2m\omega\hbar}+\sqrt{\frac{2}{m\omega\hbar}}\,\mathrm{i}p\hat a^\dagger+\frac{\left(\hat a^\dagger\right)^2}{2}\right]
    \end{equation}
    \begin{equation}
      \begin{aligned}
        G'\left(\hat{a}^\dag \right)
        &=\left(\sqrt{\frac{2}{m\omega\hbar}}\,\mathrm{i}p + \hat{a}^\dag \right) G\left(\hat{a}^\dag \right),
      \end{aligned}
    \end{equation}
    可得
    \begin{equation}
      \begin{aligned}
        \hat{p} \Ket{p }
        &=-\mathrm{i} \sqrt{\frac{\hbar m \omega }{2 } } \left(\hat{a}-\hat{a}^\dag \right) \left\{\left(\frac{1}{m\omega\pi\hbar}\right)^{1/4}\exp\left[-\frac{p^2}{2m\omega\hbar}+\sqrt{\frac{2}{m\omega\hbar}}\,\mathrm{i}p\hat a^\dagger+\frac{\left(\hat a^\dagger\right)^2}{2}\right]\Ket{0} \right\} \\
        &=-\mathrm{i} \sqrt{\frac{\hbar m \omega }{2 } } \left(\frac{1}{m\omega\pi\hbar}\right)^{1/4} \left(\hat{a}-\hat{a}^\dag \right) G\left(\hat{a}^\dag \right) \Ket{0 } \\
        &=-\mathrm{i} \sqrt{\frac{\hbar m \omega }{2 } } \left(\frac{1}{m\omega\pi\hbar}\right)^{1/4} \left\{\left[\hat{a},G\left(\hat{a}^\dag \right) \right]-\hat{a}^\dag G\left(\hat{a}^\dag \right) \right\}  \Ket{0 } \\
        &=-\mathrm{i} \sqrt{\frac{\hbar m \omega }{2 } } \left(\frac{1}{m\omega\pi\hbar}\right)^{1/4} \left\{G'\left(\hat{a}^\dag \right)-\hat{a}^\dag G\left(\hat{a}^\dag \right) \right\}  \Ket{0 } \\
        &=-\mathrm{i} \sqrt{\frac{\hbar m \omega }{2 } } \left(\frac{1}{m\omega\pi\hbar}\right)^{1/4} \left\{\left(\sqrt{\frac{2}{m\omega\hbar}}\,\mathrm{i}p + \hat{a}^\dag \right) G\left(\hat{a}^\dag \right)-\hat{a}^\dag G\left(\hat{a}^\dag \right) \right\}  \Ket{0 } \\
        &=p \left(\frac{1}{m\omega\pi\hbar}\right)^{1/4} G\left(\hat{a}^\dag \right) \Ket{0 } \\
        &=p\Ket{p }.
      \end{aligned}
    \end{equation}
  \end{enumerate}
\end{answer}

\begin{problem}
位移算符 $\hat D\left(\alpha\right)$ 定义为
\begin{equation}
  \hat D\left(\alpha\right)
  =\exp\left(\alpha\hat a^\dagger-\alpha^*\hat a\right).
\end{equation}
\begin{enumerate}
  \item 证明
  \begin{equation}
    \hat D\left(\alpha\right)\hat a\hat D^\dagger\left(\alpha\right)
    =\hat a-\alpha,\qquad
    \hat D\left(\alpha\right)\hat a^\dagger\hat D^\dagger\left(\alpha\right)
    =\hat a^\dagger-\alpha^*.
  \end{equation}
  \item 求
  \begin{equation}
    \hat D\left(\alpha\right)\hat D\left(\beta\right)\hat x^2
    \hat D^\dagger\left(\beta\right)\hat D^\dagger\left(\alpha\right).
  \end{equation}
\end{enumerate}
\end{problem}

\begin{answer}
  \begin{enumerate}
    \item Hadamard 引理给出：
    \begin{equation}
      \mathrm{e}^{A} B \mathrm{e}^{-A}
      =B + \left[A,B \right] + \frac{1 }{2! } \left[A,\left[A,B \right] \right] + \cdots.
    \end{equation}
    令
    \begin{equation}
      A = \alpha a^\dag - \alpha^*a,\qquad
      D(\alpha) = \mathrm{e}^{A},\qquad
      D^\dag(\alpha) = \mathrm{e}^{-A}.
    \end{equation}
    注意到
    \begin{equation}
      \left[A,a \right]
      =-\alpha,\qquad
      \left[A,\left[A,a \right] \right]
      =0,\qquad \cdots,
    \end{equation}
    于是有
    \begin{equation}
      D(\alpha) a D^\dag(\alpha)
      =\mathrm{e}^{A} a \mathrm{e}^{-A}
      =a - \alpha.
    \end{equation}
    注意到
    \begin{equation}
      \left[A,a^\dag \right]
      =-\alpha^*,\qquad
      \left[A,\left[A,a^\dag \right] \right]
      =0,\qquad \cdots,
    \end{equation}
    于是有
    \begin{equation}
      D(\alpha) a^\dag D^\dag(\alpha)
      =\mathrm{e}^{A} a^\dag \mathrm{e}^{-A}
      =a^\dag - \alpha^*.
    \end{equation}
    \item 由
    \begin{equation}
      \hat{x}
      =\sqrt{\frac{\hbar }{2m\omega } }\left(\hat{a}+\hat{a}^\dag \right)
    \end{equation}
    可得
    \begin{equation}
      \begin{aligned}
        \hat{D}(\alpha) \hat{x} \hat{D}^\dag(\alpha)
        &=\sqrt{\frac{\hbar }{2m\omega } } \hat{D}(\alpha) \left(\hat{a}+\hat{a}^\dag \right) \hat{D}^\dag(\alpha)
        =\sqrt{\frac{\hbar }{2m\omega } } \left(\hat{a}+\hat{a}^\dag - \alpha-\alpha^* \right) \\
        &=\hat{x} - \sqrt{\frac{\hbar }{2m\omega } } \left(\alpha+\alpha^* \right)
        \equiv \hat{x} - x_\alpha.
      \end{aligned}
    \end{equation}
    于是
    \begin{equation}
      \begin{aligned}
        \hat D\left(\alpha\right)\hat D\left(\beta\right)\hat x^2\hat D^\dagger\left(\beta\right)\hat D^\dagger\left(\alpha\right)
        &=\hat D\left(\alpha\right)\hat D\left(\beta\right)\hat{x} \hat{D}^\dag(\beta) \hat{D}(\beta) \hat{x} \hat D^\dagger\left(\beta\right)\hat D^\dagger\left(\alpha\right) \\
        &=\hat{D}(\alpha) \left(\hat{x}-x_\beta \right) \left(\hat{x}-x_\beta \right) \hat{D}^\dag(\alpha) \\
        &=\hat{D}(\alpha) \left(\hat{x}-x_\beta \right)\hat{D}^\dag(\alpha) \hat{D}(\alpha) \left(\hat{x}-x_\beta \right) \hat{D}^\dag(\alpha) \\
        &=\left(\hat{x}-x_\alpha-x_\beta \right)^2 \\
        &=\left[\hat{x}-\sqrt{\frac{\hbar }{2m\omega } }\left(\alpha+\alpha^*+\beta+\beta^* \right) \right]^2.
      \end{aligned}
    \end{equation}
  \end{enumerate}
\end{answer}

\begin{problem}
求在相干态表象下，算符 $\hat x$ 和 $\hat p$ 的外积表达式。
\end{problem}

\begin{answer}
  相干态是湮灭算符本征态：
  \begin{equation}
    \hat{a}\Ket{\alpha } = \alpha\Ket{\alpha },\qquad
    \Bra{\alpha } \hat{a}^\dag = \alpha^* \Bra{\alpha }.
  \end{equation}
  由相干态完备性关系
  \begin{equation}
    \int\frac{\mathrm{d}^2\alpha }{\pi } \Ket{\alpha }\Bra{\alpha }
    =\hat{I}, 
  \end{equation}
  可得
  \begin{equation}
    \begin{aligned}
      \hat{a}
      &=\hat{a} \hat{I}
      =\int\frac{\mathrm{d}^2\alpha }{\pi } \hat{a} \Ket{\alpha }\Bra{\alpha }
      =\int\frac{\mathrm{d}^2\alpha }{\pi } \alpha \Ket{\alpha }\Bra{\alpha },
    \end{aligned}
  \end{equation}
  \begin{equation}
    \begin{aligned}
      \hat{a}^\dag
      &=\hat{I} \hat{a}^\dag
      =\int\frac{\mathrm{d}^2\alpha }{\pi } \Ket{\alpha }\Bra{\alpha } \hat{a}^\dag 
      =\int\frac{\mathrm{d}^2\alpha }{\pi } \alpha^* \Ket{\alpha }\Bra{\alpha }.
    \end{aligned}
  \end{equation}
  因此
  \begin{equation}
    \hat{x}
    =\sqrt{\frac{\hbar }{2m\omega } }\left(\hat{a}+\hat{a}^\dag \right)
    =\sqrt{\frac{\hbar }{2m\omega } }\int\frac{\mathrm{d}^2\alpha }{\pi } \left(\alpha+\alpha^* \right) \Ket{\alpha }\Bra{\alpha },
  \end{equation}
  \begin{equation}
    \hat{p}
    =-\mathrm{i} \sqrt{\frac{\hbar m \omega }{2 } } \left(\hat{a}-\hat{a}^\dag \right)
    =-\mathrm{i} \sqrt{\frac{\hbar m \omega }{2 } } \int\frac{\mathrm{d}^2\alpha }{\pi } \left(\alpha-\alpha^* \right) \Ket{\alpha }\Bra{\alpha }.
  \end{equation}
\end{answer}

\begin{problem}
计算在相干态下，谐振子系统粒子动量的几率分布
$\left|\Braket{p|\alpha}\right|^2$。
\end{problem}

\begin{answer}
  考虑相干态本征方程：
  \begin{equation}
    \hat{a} \Ket{\alpha } = \alpha \Ket{\alpha }.
  \end{equation}
  记
  \begin{equation}
    \Braket{p|\alpha } \equiv \psi_\alpha(p),
  \end{equation}
  本征方程两边左乘 $\Bra{p } $ 得
  \begin{equation}
    \Braket{p|\hat{a}|\alpha } = \alpha\Braket{p|\alpha } = \alpha \psi_\alpha(p).
    \label{eq:eigeq-p}
  \end{equation}
  利用
  \begin{equation}
    \hat{a}
    =\sqrt{\frac{m\omega }{2\hbar } } \left(\hat{x} + \frac{\mathrm{i} }{m\omega }\hat{p}  \right),
  \end{equation}
  以及
  \begin{equation}
    \Braket{p|\hat{x}|\alpha }
    =\mathrm{i} \hbar \frac{\partial }{\partial p } \Braket{p|\alpha } 
    =\mathrm{i} \hbar \frac{\mathrm{d}}{\mathrm{d} p } \psi_\alpha(p),
  \end{equation}
  \begin{equation}
    \Braket{p|\hat{p}|\alpha }
    =p\Braket{p|\alpha }
    =p\psi_\alpha(p),
  \end{equation}
  代入\eqref{eq:eigeq-p}可得：
  \begin{equation}
    \sqrt{\frac{m\omega }{2\hbar } } \left[\mathrm{i}\hbar \frac{\mathrm{d} }{\mathrm{d}p } \psi_\alpha(p) + \frac{\mathrm{i} }{m\omega } p \psi_\alpha(p) \right]
    =\alpha \psi_\alpha(p).
  \end{equation}
  也即：
  \begin{equation}
    \frac{\mathrm{d}\psi_\alpha(p) }{\psi_\alpha(p) }
    =\left[-\frac{1 }{\hbar m\omega }p-\mathrm{i} \sqrt{\frac{2 }{\hbar m\omega } }\alpha  \right]\mathrm{d}p, 
  \end{equation}
  两边积分得到通解：
  \begin{equation}
    \psi_\alpha(p)
    =C \exp\left(-\frac{1 }{2\hbar m\omega }p^2 - \mathrm{i} \sqrt{\frac{2 }{\hbar m\omega } }\alpha p \right).
  \end{equation}
  设
  \begin{equation}
    \alpha = \alpha_{\mathrm{R}} + \mathrm{i}\alpha_{\mathrm{I}},\qquad \alpha_{\mathrm{R}},\alpha_{\mathrm{I}}\in\mathbb{R},
  \end{equation}
  则
  \begin{equation}
    \left|\psi_\alpha(p) \right|^2
    =\left|C \right|^2 \exp\left(-\frac{1 }{\hbar m\omega }p^2 + 2\sqrt{\frac{2 }{\hbar m\omega } } \alpha_{\mathrm{I}} p \right),
  \end{equation}
  由归一化条件
  \begin{equation}
    \int_{-\infty}^{+\infty} \left|\psi_\alpha(p) \right|^2 \mathrm{d}p
    =1,
  \end{equation}
  以及高斯积分公式
  \begin{equation}
    \int_{-\infty}^{+\infty} \mathrm{e}^{-Ax^2+Bx} \mathrm{d}x
    =\sqrt{\frac{\pi }{A } }\exp\left(\frac{B^2 }{4A }  \right),\qquad A>0,
  \end{equation}
  可得归一化系数满足
  \begin{equation}
    \left|C \right|^2
    =\left(\pi\hbar m\omega \right)^{-1/2} \exp\left(-2\alpha_{\mathrm{I}}^2 \right).
  \end{equation}
  因此，在相干态下，谐振子系统粒子动量的几率分布为：
  \begin{equation}
    \begin{aligned}
      \left|\Braket{p|\alpha } \right|^2
      &=\left|\psi_\alpha(p) \right|^2
      =\frac{1 }{\sqrt{\pi\hbar m \omega} } \exp\left(-\frac{1 }{\hbar m\omega }p^2 + 2\sqrt{\frac{2 }{\hbar m\omega } } \alpha_{\mathrm{I}} p - 2\alpha_{\mathrm{I}}^2 \right) \\
      &=\frac{1 }{\sqrt{\pi\hbar m \omega} } \exp\left[-\left(\frac{1 }{\sqrt{\hbar m\omega} }p-\sqrt{2}\mathrm{Im}\,\alpha  \right)^2 \right].
    \end{aligned}
  \end{equation}
\end{answer}

\begin{problem}
写出关于谐振子第一激发态 $\Ket{1}$ 的任意三种不同的相干态展开式。
\end{problem}

\begin{answer}
  由于
  \begin{equation}
    \hat{a} \Ket{n } = \sqrt{n}\Ket{n-1 },\qquad
    \hat{a}^\dag \Ket{n } = \sqrt{n+1} \Ket{n+1 },
  \end{equation}
  于是
  \begin{equation}
    \begin{aligned}
      \hat{a}^k\left(\hat{a}^\dag \right)^k\Ket{1 }
      =\sqrt{2}\sqrt{3}\cdots\sqrt{k+1} \hat{a}^k \Ket{k+1 }
      =\left(\sqrt{2}\sqrt{3}\cdots\sqrt{k+1} \right)^2 \Ket{1 }
      =(k+1)! \Ket{1 },
    \end{aligned}
  \end{equation}
  \begin{equation}
    \Ket{1 }
    =\frac{1 }{(k+1)! }\hat{a}^k \left(\hat{a}^\dag \right)^k \Ket{1 }.
  \end{equation}
  在湮灭、产生算符中间插入相干态完备性关系：
  \begin{equation}
    \begin{aligned}
      \Ket{1 }
      &=\frac{1 }{(k+1)! } \int \frac{\mathrm{d}^2\alpha }{\pi } \hat{a}^k \Ket{\alpha } \Bra{\alpha } \left(\hat{a}^\dag \right)^k \Ket{1 } \\
      &=\frac{1 }{(k+1)! } \int \frac{\mathrm{d}^2\alpha }{\pi } \alpha^k \Ket{\alpha } \Bra{\alpha } \left(\alpha^* \right)^k \Ket{1 } \\
      &=\frac{1 }{(k+1)! } \int \frac{\mathrm{d}^2\alpha }{\pi } \left|\alpha \right|^{2k} \Ket{\alpha } \Braket{\alpha|1 } \\
    \end{aligned}
  \end{equation}
  把相干态用数态展开：
  \begin{equation}
    \Ket{\alpha }
    =\mathrm{e}^{-\left|\alpha \right|^2/2} \sum_{n=0}^{\infty} \frac{\alpha^n }{\sqrt{n!} } \Ket{n }, 
  \end{equation}
  两边左乘 $\Bra{1 } $ 得：
  \begin{equation}
    \Braket{1|\alpha }
    =\alpha \mathrm{e}^{-\left|\alpha \right|^2/2},\qquad
    \Braket{\alpha|1 }
    =\Braket{1|\alpha }^*
    =\alpha^* \mathrm{e}^{-\left|\alpha \right|^2/2},
  \end{equation}
  于是
  \begin{equation}
    \begin{aligned}
      \Ket{1 }
      &=\frac{1 }{(k+1)! } \int \frac{\mathrm{d}^2\alpha }{\pi } \left|\alpha \right|^{2k} \Ket{\alpha } \Braket{\alpha|1 } \\
      &=\frac{1 }{(k+1)! } \int \frac{\mathrm{d}^2\alpha }{\pi } \left|\alpha \right|^{2k} \mathrm{e}^{-\left|\alpha \right|^2/2} \alpha^* \Ket{\alpha } .
    \end{aligned}
  \end{equation}
  分别取 $k=0,1,2 $ 就得到 $\Ket{1 } $ 在相干态表象下的三种展开式：
  \begin{equation}
    \Ket{1 }
    =\int \frac{\mathrm{d}^2\alpha }{\pi } \mathrm{e}^{-\left|\alpha \right|^2/2} \alpha^* \Ket{\alpha },
  \end{equation}
  \begin{equation}
    \Ket{1 }
    =\frac{1 }{2 } \int \frac{\mathrm{d}^2\alpha }{\pi } \left|\alpha \right|^{2} \mathrm{e}^{-\left|\alpha \right|^2/2} \alpha^* \Ket{\alpha },
  \end{equation}
  \begin{equation}
    \Ket{1 }
    =\frac{1 }{6 } \int \frac{\mathrm{d}^2\alpha }{\pi } \left|\alpha \right|^{4} \mathrm{e}^{-\left|\alpha \right|^2/2} \alpha^* \Ket{\alpha }.
  \end{equation}
\end{answer}

\begin{problem}
证明
\begin{equation}
  \int_{\mathbb{C}}\frac{\mathrm{d}^2\alpha}{\pi}\mathrm{e}^{\beta^*\alpha-\left|\alpha\right|^2}\left(\alpha^*\right)^n
  =\left(\beta^*\right)^n.
\end{equation}
\end{problem}

\begin{answer}
  令
  \begin{equation}
    I_n
    =\int\frac{\mathrm{d}^2\alpha}{\pi}\mathrm{e}^{\beta^*\alpha-\left|\alpha\right|^2}\left(\alpha^*\right)^n,
  \end{equation}
  考虑 $\mathrm{e}^{\beta^*\alpha} $ 的级数展开，有
  \begin{equation}
    \begin{aligned}
      I_n
      &=\int\frac{\mathrm{d}^2\alpha}{\pi}\mathrm{e}^{\beta^*\alpha-\left|\alpha\right|^2}\left(\alpha^*\right)^n
      =\int\frac{\mathrm{d}^2\alpha}{\pi} \left(\sum_{m=0}^{\infty} \frac{\left(\beta^* \right)^m \alpha^m }{m! }  \right) \mathrm{e}^{-\left|\alpha\right|^2}\left(\alpha^*\right)^n \\
      &=\sum_{m=0}^{\infty} \frac{\left(\beta^* \right)^m }{m! } \int \frac{\mathrm{d}^2 \alpha }{\pi } \mathrm{e}^{-\left|\alpha \right|^2} \alpha^m \left(\alpha^* \right)^n.
    \end{aligned}
  \end{equation}
  令
  \begin{equation}
    J_{mn}
    =\int \frac{\mathrm{d}^2 \alpha }{\pi } \mathrm{e}^{-\left|\alpha \right|^2} \alpha^m \left(\alpha^* \right)^n,
  \end{equation}
  考虑极坐标
  \begin{equation}
    \alpha=r\mathrm{e}^{\mathrm{i}\phi},\qquad
    \mathrm{d}^2\alpha = r\mathrm{d}r\mathrm{d}\phi,
  \end{equation}
  则
  \begin{equation}
    \begin{aligned}
      J_{mn}
      &=\int \frac{\mathrm{d}^2 \alpha }{\pi } \mathrm{e}^{-\left|\alpha \right|^2} \alpha^m \left(\alpha^* \right)^n
      =\int \frac{r\mathrm{d}r\mathrm{d}\phi }{\pi } \mathrm{e}^{-r^2} \left(r\mathrm{e}^{\mathrm{i}\phi} \right)^m \left(r\mathrm{e}^{-\mathrm{i}\phi} \right)^n \\
      &=\frac{1 }{\pi } \int_{0}^{\infty}\mathrm{d}r\, r^{m+n+1} \mathrm{e}^{-r^2} \int_{0}^{2\pi}\mathrm{d}\phi\, \mathrm{e}^{\mathrm{i}\left(m-n \right)\phi}
      =\frac{1 }{\pi } \int_{0}^{\infty}\mathrm{d}r\, r^{m+n+1} \mathrm{e}^{-r^2} \cdot 2\pi \delta_{mn} \\
      &=2\delta_{mn} \int_{0}^{\infty}\mathrm{d}r\, r^{2n+1} \mathrm{e}^{-r^2}.
    \end{aligned}
  \end{equation}
  令
  \begin{equation}
    u=r^2,\qquad
    \mathrm{d}u=2r\mathrm{d}r,
  \end{equation}
  则
  \begin{equation}
    \begin{aligned}
      J_{mn}
      =2\delta_{mn} \int_{0}^{\infty}\mathrm{d}r\, r^{2n+1} \mathrm{e}^{-r^2}
      =\delta_{mn} \int_{0}^{\infty}\mathrm{d}u\, u^n \mathrm{e}^{-u}
      =\delta_{mn} \Gamma(n+1)
      =\delta_{mn} n!.
      \label{eq:J_mn}
    \end{aligned}
  \end{equation}
  于是
  \begin{equation}
    \begin{aligned}
      I_n
      &=\sum_{m=0}^{\infty} \frac{\left(\beta^* \right)^m }{m! } \int \frac{\mathrm{d}^2 \alpha }{\pi } \mathrm{e}^{-\left|\alpha \right|^2} \alpha^m \left(\alpha^* \right)^n
      =\sum_{m=0}^{\infty} \frac{\left(\beta^* \right)^m }{m! } J_{mn} \\
      &=\sum_{m=0}^{\infty} \frac{\left(\beta^* \right)^m }{m! } \delta_{mn} n!
      =\left(\beta^* \right)^n.
    \end{aligned}
  \end{equation}
\end{answer}

\begin{problem}
通过计算表达式的右方，验证谐振子粒子数算符的本征态 $\Ket{n}$ 的如下展开式：
\begin{enumerate}
  \item
  \begin{equation}
    \Ket{n}=\int_{\mathbb{C}}\frac{\mathrm{d}^2\alpha}{\pi}\Ket{\alpha}\exp\left(-\left|\alpha\right|^2/2\right)
    \frac{\left(\alpha^*\right)^n}{\sqrt{n!}}.
  \end{equation}
  \item
  \begin{equation}
    \Ket{n}=\frac{1}{2\pi}\exp\left(\frac{r_0^2}{2}\right)r_0^{-n}\sqrt{n!}\int_0^{2\pi}\mathrm{e}^{-\mathrm{i}n\phi}\Ket{r_0\mathrm{e}^{\mathrm{i}\phi}}\mathrm{d}\phi,
  \end{equation}
  其中 $\alpha\equiv r_0\mathrm{e}^{\mathrm{i}\phi}$。
\end{enumerate}
\end{problem}

\begin{answer}
  \begin{enumerate}
    \item 把相干态用数态展开，再利用 \eqref{eq:J_mn}，有
    \begin{equation}
      \begin{aligned}
        &\int\frac{\mathrm{d}^2\alpha}{\pi}\Ket{\alpha}\exp\left(-\left|\alpha\right|^2/2\right)\frac{\left(\alpha^*\right)^n}{\sqrt{n!}} \\
        =&\int\frac{\mathrm{d}^2\alpha}{\pi}\left(\mathrm{e}^{-\left|\alpha \right|^2/2} \sum_{m=0}^{\infty} \frac{\alpha^m }{\sqrt{m!} } \Ket{m }  \right)\exp\left(-\left|\alpha\right|^2/2\right)\frac{\left(\alpha^*\right)^n}{\sqrt{n!}} \\
        =&\sum_{m=0}^{\infty} \frac{1 }{\sqrt{m!n!} } \int \frac{\mathrm{d}^2\alpha }{\pi } \mathrm{e}^{-\left|\alpha \right|^2} \alpha^m \left(\alpha^* \right)^n \Ket{m }
        =\sum_{m=0}^{\infty} \frac{1 }{\sqrt{m!n!} } J_{mn} \Ket{m } \\
        =&\sum_{m=0}^{\infty} \frac{1 }{\sqrt{m!n!} } \left(\delta_{mn} n! \right) \Ket{m } \\
        =&\Ket{n }.
      \end{aligned}
    \end{equation}
    \item 相干态按数态展开：
    \begin{equation}
      \Ket{\alpha }
      =\Ket{r_0\mathrm{e}^{\mathrm{i}\phi} }
      =\left(\mathrm{e}^{-\left|\alpha \right|^2/2} \sum_{m=0}^{\infty} \frac{\alpha^m }{\sqrt{m!} } \Ket{m }  \right)
      =\mathrm{e}^{-r_0^2/2} \sum_{m=0}^{\infty} \frac{r_0^m \mathrm{e}^{\mathrm{i}m\phi} }{\sqrt{m!} } \Ket{m },
    \end{equation}
    因此
    \begin{equation}
      \begin{aligned}
        &\frac{1}{2\pi}\exp\left(\frac{r_0^2}{2}\right)r_0^{-n}\sqrt{n!}\int_0^{2\pi}\mathrm{e}^{-\mathrm{i}n\phi}\Ket{r_0\mathrm{e}^{\mathrm{i}\phi}}\mathrm{d}\phi \\
        =&\frac{1}{2\pi}\exp\left(\frac{r_0^2}{2}\right)r_0^{-n}\sqrt{n!}\int_0^{2\pi}\mathrm{e}^{-\mathrm{i}n\phi} \left(\mathrm{e}^{-r_0^2/2} \sum_{m=0}^{\infty} \frac{r_0^m \mathrm{e}^{\mathrm{i}m\phi} }{\sqrt{m!} } \Ket{m } \right) \mathrm{d}\phi \\
        =&\frac{1}{2\pi}r_0^{-n}\sqrt{n!}\sum_{m=0}^{\infty} \frac{r_0^m }{\sqrt{m!} } \Ket{m } \int_0^{2\pi}\mathrm{e}^{\mathrm{i}\left(m-n \right)\phi}  \mathrm{d}\phi \\
        =&\frac{1}{2\pi}r_0^{-n}\sqrt{n!}\sum_{m=0}^{\infty} \frac{r_0^m }{\sqrt{m!} } \Ket{m } \cdot 2\pi\delta_{mn} \\
        =&\Ket{n }.
      \end{aligned}
    \end{equation}
  \end{enumerate}
\end{answer}

\begin{problem}
设粒子在 $\delta$ 势阱中运动，Hamilton 算符为
\begin{equation}
  \hat H=\frac{\hat p^2}{2m}-g\delta\left(\hat x\right).
\end{equation}
试在动量表象中求解束缚能级 $E<0$ 和对应本征态。
\end{problem}

\begin{answer}
  设\begin{equation}
    \hat{H} \Ket{\psi }
    =\left[\frac{\hat p^2}{2m}-g\delta\left(\hat x\right) \right]\Ket{\psi }
    =E \Ket{\psi },
  \end{equation}
  两边左乘 $\Bra{p } $ 得：
  \begin{equation}
    \frac{p^2 }{2m } \Braket{p|\psi } - g\Braket{p|\delta\left(\hat{x} \right)|\psi }
    =E\Braket{p|\psi }. 
  \end{equation}
  记
  \begin{equation}
    \psi(p) = \Braket{p|\psi },
  \end{equation}
  注意到
  \begin{equation}
    \begin{aligned}
      \Braket{p|\delta\left(\hat{x} \right)|\psi }
      &=\int \mathrm{d}p' \Braket{p|\delta\left(\hat{x} \right)|p' }\Braket{p'|\psi }
      =\int \mathrm{d}p' \int\mathrm{d}x \Braket{p|\delta\left(\hat{x} \right)|x }\Braket{x|p' } \psi(p') \\
      &=\int \mathrm{d}p' \int\mathrm{d}x \delta\left(x \right) \Braket{p|x }\Braket{x|p' } \psi(p') \\
      &=\int \mathrm{d}p' \Braket{p|x=0 }\Braket{x=0|p' } \psi(p') \\
      &=\frac{1 }{2\pi\hbar } \int_{-\infty}^{+\infty}\mathrm{d}p' \psi(p')
      \equiv \frac{I }{2\pi\hbar },\qquad
      I=\int_{-\infty}^{+\infty}\mathrm{d}p' \psi(p')
    \end{aligned}
  \end{equation}
  因此有
  \begin{equation}
    \psi(p)
    =\frac{gI }{2\pi\hbar\left(\frac{p^2 }{2m }-E  \right) }, 
  \end{equation}
  两边对 $p $ 积分：
  \begin{equation}
    I
    =\frac{gI }{2\pi\hbar } \int_{-\infty}^{+\infty} \frac{1 }{\frac{p^2 }{2m }-E  } \mathrm{d}p
    =\frac{gI }{2\pi\hbar } \cdot \pi\sqrt{\frac{2m }{-E } },
  \end{equation}
  对比头尾可得本征能量：
  \begin{equation}
    E
    =-\frac{mg^2 }{2\hbar^2 } 
  \end{equation}
  设动量表象本征波函数归一化为
  \begin{equation}
    \psi(p)
    =\frac{A }{\left(\frac{p^2 }{2m }-E  \right) },
  \end{equation}
  由归一化条件
  \begin{equation}
    \int_{-\infty}^{+\infty} \left|\psi(p) \right|^2\mathrm{d}p
    =1
  \end{equation}
  可得 $A $ 可取为
  \begin{equation}
    A
    =\sqrt{\frac{m }{2\pi } } \left(\frac{g }{\hbar }  \right)^{3/2},
  \end{equation}
  于是
  \begin{equation}
    \psi(p)
    =\sqrt{\frac{m }{2\pi } } \left(\frac{g }{\hbar }  \right)^{3/2} \frac{1 }{\frac{p^2 }{2m } + \frac{mg^2 }{2\hbar^2 }   }.
  \end{equation}
  态矢量为
  \begin{equation}
    \Ket{\psi }
    =\int \mathrm{d}p \psi(p) \Ket{p }
    =\sqrt{\frac{m }{2\pi } } \left(\frac{g }{\hbar }  \right)^{3/2} \int \mathrm{d}p \frac{1 }{\frac{p^2 }{2m } + \frac{mg^2 }{2\hbar^2 }   } \Ket{p }
    =\sqrt{\frac{2 }{\pi } } \left(\frac{mg }{\hbar }  \right)^{3/2} \int \mathrm{d}p \frac{\Ket{p } }{p^2+\left(gm/\hbar \right)^2 }. 
  \end{equation}
\end{answer}

\end{document}
